% QMB_n15_zeit_superposition.tex — Kapitel 15: Zeit im Zustandsraum, Superposition ohne Zeit
% Inhaltliche Grundlage: Dok. 307 (Juli 2026), Dok. 334 (Aug. 2026), Dok. 333; eigene Darstellung
\chapter{Wo die Zeit sitzt: Zustandsraum und Superposition}
\label{ch:zeit}

\section{Die Gabelung}

Am Anfang jeder quantenmechanischen Modellbildung steht eine Wahl, die
meist unausgesprochen bleibt und dennoch alles Weitere festlegt: Ist
die Zeit eine Richtung \emph{im} Zustandsraum, oder läuft sie
\emph{daneben} her, als äußerer Parameter?

Die Frage klingt technisch, ist aber ontologisch. Sie entscheidet nicht,
welche Mathematik verwendet wird --- beide Antworten arbeiten mit
selbstadjungierten Operatoren, Spektren und Projektoren ---, sondern
\emph{wo} die Zeit sitzt. Und genau diese Verortung entscheidet, ob die
physikalischen Größen aus der Struktur \emph{fallen} oder ihr
\emph{angeheftet} werden müssen.

Die Frage reicht zu den Wurzeln der Quantentheorie zurück. Schon 1932
erkannte Wolfgang Pauli, dass ein Zeitoperator, der kanonisch zum
Hamiltonoperator konjugiert ist, auf grundsätzliche Schwierigkeiten
stößt. Das Pauli-Theorem hat die Entwicklung der Quantenmechanik
geprägt und die Zeit als äußeren Parameter scheinbar unumstößlich
gemacht. Spätere Arbeiten --- Aharonov und Bohm, Garrison und Wong ---
zeigten, dass unter bestimmten Bedingungen, etwa bei nicht
halbbeschränkten Hamiltonoperatoren oder kompakten Konfigurationen,
Zeitoperatoren doch konstruierbar sind. Diese Ausnahme ist für die FFGFT
der Regelfall.

\section{Die ausgerollte Sicht: Zeit als äußerer Parameter}

Die vertraute Formulierung siedelt die Zeit außerhalb des Zustandsraums
an. Der Hilbertraum ist der Raum der Zustände \emph{zu einem
Zeitpunkt}, typischerweise $\mathcal{H} = L^2(\text{Konfigurationsraum})$.
Die Zeit ist der Parameter $t$, entlang dessen sich der Zustand
entwickelt, vermittelt durch $e^{-i\hat Ht}$. Ort, Impuls, Spin sind
Operatoren auf $\mathcal{H}$. Die Zeit ist keiner.

Dass sie keiner sein kann, ist kein Versehen, sondern ein Resultat:
Paulis Einwand zeigt, dass zu einem nach unten beschränkten
Hamiltonoperator kein selbstadjungierter, kanonisch konjugierter
Zeitoperator existiert.\status{E} In der ausgerollten Sicht bleibt die
Zeit deshalb notwendig die Bühne, nicht das Stück.

Die Konsequenz ist strukturell: Da die Zeit außen steht, enthält der
Zustandsraum allein keine dynamischen Skalen. Massen, Kopplungen,
Lebensdauern sind in der reinen Raumstruktur nicht vorhanden --- sie
müssen von außen eingeführt werden. Im Standardmodell geschieht das
durch Yukawa-Terme und die freien Parameter, die deren Stärke
festlegen. Die Struktur liefert den Rahmen; die Zahlen kommen aus der
Messung.

Das Problem ist alt. Weyl und Wigner erkannten, dass die Massen der
Elementarteilchen gruppentheoretisch als Casimir-Invarianten erscheinen
--- als strukturelle Label, nicht als angeheftete Parameter. Dennoch
blieb die Herkunft der Massenwerte offen. Der Higgs-Mechanismus erklärt
die Masse der Eichbosonen, führt aber selbst neue freie Parameter ein.
Die Kritik an der ausgerollten Sicht ist nicht, dass sie falsch wäre,
sondern dass sie die Herkunft der Massenskalen unbeantwortet lässt und
auf externe Eingaben verweist.

\subsection*{Eine repräsentative Konstruktion dieser Art}

Viele Modellprogramme, die eine geometrische Herkunft der
Teilcheneigenschaften anstreben, folgen bei aller Verschiedenheit
demselben Schema: Man wählt einen diskreten Träger (ein Gitter, eine
Cut-and-Project-Struktur, einen periodischen Graphen); definiert darauf
eine selbstadjungierte Operatorfamilie; extrahiert über
Spektralprojektoren Sektoren; und ordnet diesen Sektoren nachträglich
physikalische Bedeutung zu --- Ladung, Spin, Masse, Letztere meist über
eine parametrische Abbildung mit ein oder zwei freien Parametern.

Ein solches Vorgehen kann mathematisch vollständig kontrolliert sein.
Die entscheidende Beobachtung betrifft nicht seine Strenge, sondern
seine Verortung der Zeit: Der Träger ist ein räumliches Objekt, die
Operatorevolution läuft über einen äußeren Parameter. Die Sektoren sind
räumliche Eigenräume, keine Zeitmoden. Die Masse ist in ihnen nicht
enthalten; sie wird angeheftet.

Und hier greift ein Argument, das sich selbst trägt: \textbf{Eine
parametrische Abbildung mit freien Parametern kann jede endliche
Werteliste treffen.} Ist die Abbildung nicht vorab und universell
fixiert, besitzt sie keinen prädiktiven Gehalt --- sie benennt die
bekannten Eingaben um. Um Zahlen zu liefern, muss eine solche
Konstruktion an externe Werte gekoppelt werden: an die Messung, und
damit indirekt an das Standardmodell selbst. Das ist redlich, sofern es
seine Grenze deklariert: dass es keine Massen ableitet, sondern eine
Hypothesenarchitektur bereitstellt. Die Schwierigkeit dieses Wegs ist
kein Zufall, sondern die Folge der ausgerollten Verortung der Zeit.

Ein konkretes Beispiel: Die Graphen-Modellierung definiert
Quasiteilchen auf hexagonalen Gittern. Die Anregungen zeigen lineare
Dispersion und lassen sich als masselose Fermionen lesen --- aber die
Einführung einer Masse erfordert stets zusätzliche Terme (Haldane-Term,
Kekulé-Distorsion) mit freien Parametern. Die Geometrie liefert die
Kinematik, nicht die Massenskala.

\section{Die eingerollte Sicht: Zeit als kompakte Dimension}

Die Alternative belässt die Zeit nicht außen und eliminiert sie auch
nicht, sondern zieht sie \emph{in} den Zustandsraum --- als geschlossene,
periodische Dimension. An die Stelle von $L^2(\text{Raum})$ mit äußerem
$t$ tritt ein Raum, in dem eine kompakte Zeitrichtung Teil der Geometrie
ist: $\mathcal{H} = L^2(T^4)\otimes\mathbb{C}^3$, wobei eine der
Torusrichtungen die Zeit trägt. Die Bijektivität dieser Darstellung
(Typ III) ist im Korpus belegt (Kapitel~\ref{ch:hilbert}).

Auf einer geschlossenen Dimension sind die zulässigen Zustände stehende
Moden mit diskreten Windungszahlen. Die Quantisierung ist dann nicht
postuliert, sondern geometrisch erzwungen: Nur bestimmte Moden passen
auf die geschlossene Struktur. Die Masse wird zum Eigenwert eines
Operators auf $\mathcal{H}$ selbst --- nicht angeheftet, sondern
strukturell enthalten. Die Verhältnisse der Massen liegen fest, weil die
zulässigen Moden festliegen.\status{B}

Die Relation $\tilde T\cdot m = 1$ ist die konkrete Ausführung dieses
Gedankens: Jede massive Größe \emph{ist} ihr eigener Zeitzyklus. Die
Entkompaktifizierung $S^1_m\to\mathbb{R}_t$ ist maßerhaltend, also
verlustfrei. Dass eine kompakte Dimension einen diskreten Modenturm
erzeugt, dessen Massen sich wie das Inverse der Kompaktifizierungsperiode
verhalten, ist kein exotischer Zug, sondern der Mechanismus, den Kaluza
und Klein für eine kompakte Raumdimension vorgeführt haben --- hier auf
die Zeit angewandt.

Paulis Einwand trifft die eingerollte Sicht nicht. Er gilt für eine
offene, nach unten unbeschränkte Zeit konjugiert zu einem beschränkten
Hamiltonoperator. Eine kompakte Zeit verhält sich wie eine
Winkelvariable: Ihr konjugiertes Spektrum ist diskret, wie beim
Drehimpuls zum Winkel, und das Pauli-Argument entfällt.\status{B} Die
geschlossene Zeit ist nicht nur zulässig, sie ist der natürliche Ort
einer diskreten Modenstruktur.

Diese Einsicht hat Wurzeln: Witten wies darauf hin, dass kompakte
Zeitdimensionen in der Stringtheorie auftreten; Hawking diskutierte die
Implikationen geschlossener Zeit für die Kosmologie; Bars entwickelte
die zweizeitige Physik. Und die eingerollte Sicht berührt die
thermodynamische Zeit: Ist die Zeit kompakt, ist die Pfeilrichtung nicht
fundamental, sondern eine emergente Eigenschaft der nichtkompakten
Raumanteile --- eine natürliche Erklärung makroskopischer
Irreversibilität ohne Annahmen über Anfangsbedingungen.

\begin{laierspur}
Man kann eine Uhr auf zwei Weisen bauen. Entweder man hängt ein
Pendel an einen Rahmen und lässt die Zeit von außen laufen --- dann
muss man die Länge des Pendels von Hand einstellen. Oder man baut die
Uhr als geschlossenen Kreis, in dem jede Schwingung wieder in sich
zurückläuft --- dann bestimmt der Kreis selbst, welche Schwingungen
möglich sind. Die FFGFT baut die Uhr auf die zweite Art. Deshalb muss
sie die Massen nicht einstellen: Sie fallen aus der Geometrie.
\end{laierspur}

\section{Der Kern der Abwägung}

\textbf{Was die ausgerollte Sicht gewinnt:} den nahtlosen Anschluss an
die Maschinerie der Quantenfeldtheorie; die Freiheit von einer starken
ontologischen Festlegung; die Vertrautheit der offenen Zeitachse.
\textbf{Was sie kostet:} Die Struktur allein legt keine Massen fest.
Dynamische Werte müssen von außen kommen; Vorhersagen hängen an
externer Kalibrierung; ein Modell, das sich als eigenständig versteht,
koppelt für seine Zahlen doch wieder an die Messung.

\textbf{Was die eingerollte Sicht gewinnt:} Die Moden und ihre
Verhältnisse fallen aus der Struktur, ohne parametrisches Anheften; die
diskrete Quantisierung ist geschenkt statt postuliert. \textbf{Was sie
kostet:} eine starke ontologische Behauptung --- Zeit als geschlossene
Dimension ---, die der gewohnten Intuition der laufenden Zeit
widerspricht. Und sie hebt nicht jede äußere Anbindung auf: Die absolute
Skala braucht weiterhin eine Einheitenbrücke, die SI-Umrechnung über
feste Konversionsfaktoren.

Der entscheidende Befund ist aber dieser: \textbf{Beide Wege benutzen
denselben Operatorformalismus.} Der Unterschied liegt nicht in der
Mathematik, sondern allein darin, ob die Zeit innen oder außen sitzt.
Man kann dieselbe Konstruktion --- selbstadjungierter Operator,
Spektralsektoren --- beibehalten und nur die Zeit in den Raum ziehen;
dann werden die Sektoren zu Moden, und die Masse, die zuvor angeheftet
werden musste, wird zum Eigenwert der Struktur. Die Wahl der
Zeitverortung ist eine Rahmenwahl, die der Modellbildung vorausgeht und
deren interpretative Reichweite festlegt.

\section{Wer die Zeit schon einmal nach innen geholt hat}

Der Gedanke, die Zeit nicht als äußeren Parameter zu behandeln, ist
keine Neuheit, sondern eine lange verfolgte Linie. Es ist redlich, die
eingerollte Sicht in dieser Linie zu verorten und ihre Besonderheit
abzugrenzen.

\textbf{Relationale Zeit.} Page und Wootters lassen die Zeit innerhalb
des Zustandsraums entstehen: als Verschränkung zwischen einem
Uhr-Subsystem und dem Rest; der Gesamtzustand ist stationär,
\enquote{Evolution ohne Evolution}. Moreva et al. und Marletto und
Vedral haben den Mechanismus präzisiert und experimentell illustriert;
Gambini und Pullin entwickelten die Montevideo-Interpretation.

\textbf{Zeitlose Quantengravitation.} Wheeler--DeWitt: Die universelle
Wellenfunktion ist zeitlos, $\hat H\Psi = 0$; das daraus folgende
\enquote{Problem der Zeit} ist die Feststellung, dass es keine äußere
Zeit gibt. Kuchař und Isham lieferten die Übersichten. Rovelli:
relationale Quantenmechanik, mit Connes die thermische-Zeit-Hypothese.
Barbour: die konsequente Elimination der Zeit zugunsten zeitloser
Konfigurationen.

\textbf{Kompakte Zeit in der Stringtheorie.} Bars und Kollegen mit
zweizeitiger Physik und erweiterter Symmetriegruppe.

\textbf{Die Modenseite.} Wigner: Masse ist eine Casimir-Invariante, ein
Label irreduzibler Darstellungen der Poincaré-Gruppe --- bereits ein
struktureller Eigenwert. Kaluza und Klein: Eine kompaktifizierte
Dimension erzeugt einen diskreten Modenturm mit gequantelten Ladungen
--- die Blaupause dafür, dass eine eingerollte Dimension Moden
hervorbringt.

Die Abgrenzung ist präzise: Der überwiegende Teil dieser Tradition
macht die Zeit relational, emergent oder zeitlos --- sie nimmt die Zeit
als Fundamentales \emph{heraus}. Die eingerollte Sicht nimmt die Zeit
\emph{hinein}: als kompakte geometrische Dimension, die Moden trägt.
Sie teilt mit Page--Wootters und Rovelli die Ablehnung der äußeren
Parameterzeit, unterscheidet sich aber im positiven Vorschlag.

\section{Mathematische Ausführung}

Der Hilbertraum der eingerollten Sicht ist
$\mathcal{H} = L^2(T^1\times M)\otimes\mathcal{H}_{\mathrm{int}}$, mit
einer kompakten Zeitdimension $T^1$ der Länge $T$ und dem gewöhnlichen
Raum $M$. Eine Funktion darauf entwickelt sich in eine Fourierreihe:
\begin{equation}
  \Psi(t,\mathbf{x}) = \sum_{n\in\mathbb{Z}}e^{i n t/T}\,\psi_n(\mathbf{x}).
\end{equation}
Die diskreten Frequenzen $\omega_n = 2\pi n/T$ entsprechen den
Massenwerten in natürlichen Einheiten. Der Operator, der sie erzeugt,
ist $\hat E = -i\,\partial_t$ mit dem Spektrum $\{2\pi n/T\}$. Die
Quantisierung der Massen ist eine direkte Konsequenz der Periodizität
der Zeit; nichts muss postuliert werden.

\textbf{Yukawa-Kopplung als Auswahlregel.} Die Kopplung zwischen Moden
ist ein Matrixelement des Wechselwirkungsoperators:
\begin{equation}
  g_{ijk}\propto\int_{T^1}e^{i(n_i-n_j-n_k)t/T}\,dt
  = T\cdot\delta_{n_i,\,n_j+n_k}.
\end{equation}
Das ist eine exakte Erhaltung des Modenindex in der Zeitrichtung. Die
Yukawa-Kopplungen des Standardmodells erscheinen als Auswahlregeln für
die Modenindizes, nicht als freie Parameter.\status{B}

\textbf{Ladungsquantisierung.} Die kompakte Zeitrichtung verhält sich
wie eine fünfte Dimension im Sinne von Kaluza--Klein, aber zeitartig.
Die Ladung ist $Q\propto n/R_T$ mit dem Radius $R_T$ der kompakten Zeit
--- Ladungsquantisierung als Konsequenz der Periodizität.

\textbf{Massenhierarchie.} Der Korpus beschreibt die Leptonmassen in
drei ineinander übersetzbaren Darstellungen: erstens grundlegend über
Quantenzahlen $(n_i,l_i,j_i)$ mit $\xi_i = \xi_0\,f(n_i,l_i,j_i)$ und
$E_i = 1/\xi_i$, wobei die Absolutwerte die fraktale Korrektur
$K_{\mathrm{frak}}$ tragen; zweitens in Leiterform
$m_i = r_i\,\xi^{\pi_i}\,v$ mit rationalen $r_i$ und der
Generationshierarchie $\pi = \frac32, 1, \frac23$; drittens als
$\mathbb{Z}_3$-Zirkulant, $\sqrt{m_k} = M\bigl(1+\sqrt2\cos(\frac29+2\pi k/3)\bigr)$,
mit der Koide-Relation $Q = 2/3$ und dem Winkel $\theta = 2/9$ als
Ausgabe der Diagonalisierung. Die dimensionslosen Verhältnisse sind
exakt; Abweichungen entstehen erst beim Übergang zu absoluten
SI-Werten, aus der SI-Umrechnung, den Brückenkonstanten $v$ und $E_0$,
der fraktalen Korrektur und der endlichen Messgenauigkeit. Das ist ein
Strukturbeweis: Die Verhältnisse folgen aus der Geometrie, die
Absolutwerte aus der Brücke.\status{K}

\section{Kritische Würdigung}

\textbf{Ontologischer Status.} Die kompakte Zeit ist im Korpus real und
geometrisch, nicht bloß Rechentechnik. Der Rand $\partial T^4 = \emptyset$
macht die Frage nach einem Anfangszeitpunkt nicht falsch beantwortet,
sondern falsch gestellt. Die Zeit ist diskret: Die FFGFT besitzt ein
minimales, von $\xi$ festgelegtes Zeitinkrement $d_1 = 100\xi_0 = 1/75$,
und die zulässigen Zustände sind diskrete Moden. Wheelers
\enquote{Zeit als Illusion} teilt die eingerollte Sicht nicht: Sie
nimmt die Zeit als real, aber geometrisch anders als gewohnt.

\textbf{Die absolute Skala: fixiert, nicht frei.} Die Massenverhältnisse
folgen aus der Modenstruktur und sind exakt. Die absolute Skala ist
kein freier Parameter: Das Verhältnisgerüst ruht auf $\xi_0 = 4/30000$,
die Absolutwerte tragen ihre SI-Korrektur über feste Brückenkonstanten
--- $v$ in den Massen, $E_0 = 7{,}398$\,MeV in $\alpha = \xi E_0^2$,
feste Zahlenfaktoren in $G$. Die Planck-Größen sind kein Eingang,
sondern folgen aus $\xi$: $\xi\to\{m_e,\dots\}\to G\to\ell_P$, und
$\hbar$ selbst ist geometrisch. Die einzige exakte $\xi$-Potenzrelation
ist $L_0/\ell_P = \xi_0$, $L_0\approx2{,}15\times10^{-39}$\,m. Der
Zeitzyklus ist nicht an die Planck-Skala gekoppelt --- die Kausalrichtung
ist umgekehrt. Es bleibt kein zu eichender Modul offen; die
Dilaton-Analogie zur Stringtheorie trägt hier nicht.

\textbf{Kausalität.} Erzeugt eine periodische Zeit geschlossene
zeitartige Kurven? Der Träger ist ein Torus mit $\pi_1\neq0$ und stabilen
Windungen, kein einfacher Kreis; die relevante Verbindung ist die
topologische Verbindung auf $T^4$ (Kapitel~\ref{ch:verschraenkung}).
Die klassische CTC-Problematik der Gödel-Lösungen betrifft den
dekompaktifizierten Lorentz-Sektor; die kompakte Masse-Zeit-Richtung ist
spektral, nicht kausal geschlossen.

\textbf{Experimentelle Überprüfung.} Die eingerollte Sicht sagt keinen
verborgenen schweren Modenturm oberhalb der Beschleunigerreichweite
voraus. Das diskrete Spektrum der kompakten Zeit \emph{ist} über
$\tilde T\cdot m = 1$ bereits das beobachtete Teilchenspektrum; es gibt
keinen zusätzlichen Kaluza-Klein-Turm. Prüfbar sind: die Teilchenmassen
selbst als $\xi$-Ableitungen (Leptonstruktur, Koide auf $10^{-5}$,
$g-2$); die Modenindex-Erhaltung als Kopplungsauswahlregel; die
CHSH-Amplitude der Ordnung $\xi$; die fraktale Dimension $D_f$ und die
Fundamentallänge $L_0$. Für eine saubere Falsifikation reichen
Verhältnistests allein nicht --- belastbar ist erst eine vorwärts
festgelegte, unabhängig prüfbare Vorhersage.

\section{Die verborgene Inkonsistenz der Standard-QM}

Die Instantanität ist kein QM-spezifisches Problem, sondern ein Erbe der
klassischen Physik. In der Newtonschen Mechanik wirkt Gravitation
instantan; das Coulomb-Potential ist instantan; Lagrange- und
Hamilton-Mechanik beschreiben den Zustand $(q(t),p(t))$ zu einem fixen
$t$ --- die Berechnung ist eine zeitlose algebraische Auswertung. Niemand
fragt, wie lange diese Berechnung dauert, weil Berechnungen in dieser
Sprache keine Dauer haben. Pauli hat nur explizit gemacht, was klassisch
stillschweigend war.

Die QM erbt das und fügt eine Spannung hinzu. \textbf{Der Zustand
$\psi(t)$ ist instantan}: eine zeitlose Größe zu einem fixen $t$, keine
Dauer, kein Prozess. \textbf{Die Auswertung $|\psi(t)|^2$ braucht
Zeit}: Um eine Wahrscheinlichkeitsdichte zu verifizieren, braucht man
viele Messungen, also echte physikalische Zeit. Der Formalismus ist
instantan, seine Verifikation ist zeitlich ausgedehnt. Die Zeit, die
für die probabilistische Auswertung gebraucht wird, steht außerhalb des
Formalismus --- sie ist weder Operator noch Variable, sondern das, was
der Experimentator aufwendet, um den Zustand überhaupt zu kennen.

Superposition $c_1|\psi_1\rangle+c_2|\psi_2\rangle$ ist damit doppelt
zeitlos: als Zustand (kein Zeitoperator, kein intrinsisches
\enquote{Wann}) und als Aussage (probabilistisch, braucht Zeit zur
Verifikation, die außerhalb liegt). Was als überlagert existiert, kann
nur durch wiederholte Messungen über echte Zeit belegt werden --- und
jede Messung zerstört die Superposition.

\section{Drei Konsequenzen, die selten gesagt werden}

In den Lehrbüchern wird Zeit als Parameter eingeführt, ohne zu betonen,
was das Fehlen eines Zeitoperators bedeutet. Es bedeutet konkret:

\textbf{Superposition ist zeitlos.} Ein Elektron im Zustand
$c_\uparrow|\uparrow\rangle+c_\downarrow|\downarrow\rangle$ hat keine
intrinsische Zeitskala. Die Frage \enquote{wie lange ist es in dieser
Superposition?} ist im Formalismus nicht stellbar. Was existiert, ist
die Kohärenzzeit $T_2$ --- eine phänomenologische Größe der
Umgebungskopplung, keine intrinsische Eigenschaft des Zustands. In der
FFGFT trägt jede Mode $\tilde T_i = 1/m_i$ als geometrische Größe. Ein
Zustand \emph{ist} seine Zeitskala --- nicht als Dauer einer
Superposition, sondern als strukturelle Eigenschaft der Mode.

\textbf{Verschränkung ohne Gleichzeitigkeit.} Zwei verschränkte
Teilchen teilen einen Zustand in $\mathcal{H}_1\otimes\mathcal{H}_2$
ohne eingebaute Zeitstruktur. Die Frage nach der Gleichzeitigkeit
einer Messung an beiden ist keine QM-Frage, sondern eine klassische
über den externen Parameter --- und damit eine relativistische. Das ist
der tiefste Grund, warum die QM mit Lorentz-Invarianz vereinbar ist:
Gleichzeitigkeit steht außerhalb des Formalismus. In der FFGFT: Wären
zwei Moden mit $m_1\neq m_2$ verschränkt, hätten sie verschiedene
$\tilde T_1\neq\tilde T_2$. Modenübergreifende Verschränkung erzeugt
inkompatible Zeitskalen --- strukturell problematischer als in der QM,
wo der Konflikt nicht auftritt, weil $t$ außen steht.

\textbf{Kollaps ohne Dauer.} Der Übergang $|\psi\rangle\to|\psi_k\rangle$
hat im Formalismus keine Dauer; er ist instantan in $t$. Wie lange der
Kollaps dauert, liegt außerhalb der Standard-QM. Ohne Zeitoperator gibt
es keine Observable \enquote{Kollapsdauer}. In der FFGFT ist der
Typ-I-Schritt (Messakt) ein geometrischer Übergang: Der Massenkreis
$S^1_m$ wird zur Zeitachse $\mathbb{R}_t$, die Windungszahl wird
verworfen. Das ist ein Informationsverlust mit natürlicher Richtung.

\section{Was $\tilde T\cdot m = 1$ strukturell ändert}

Die Grundrelation setzt Zeit nicht als externen Parameter, sondern als
geometrische Größe jeder Mode: Die Periode $R_m = \hbar/(mc^2)$ ist
intrinsisch --- sie gehört zum Teilchen wie seine Masse. Drei
Konsequenzen:

\textbf{Superposition} innerhalb einer Mode --- intermediäre Torsion
--- hat eine natürliche Zeitskala $\tilde T_i$. Superposition
\emph{über} verschiedene Moden würde inkompatible Zeitskalen
überlagern; das ist der strukturelle Grund, warum modenübergreifende
Superpositionen in der FFGFT problematisch sind.\status{B}

\textbf{Verschränkung} innerhalb einer Mode hat Gleichzeitigkeit durch
$\tilde T$ definiert. Über Moden hinweg fehlt ein gemeinsamer
Zeitbegriff --- ähnlich wie in der QM, aber aus anderem Grund: nicht
weil $t$ extern steht, sondern weil verschiedene Moden verschiedene
geometrische Zeitskalen tragen.

\textbf{Kollaps} (Typ I) hat eine natürliche Richtung ---
verlustbehaftet von $S^1_m$ nach $\mathbb{R}_t$ --- und eine natürliche
Asymmetrie: Die Rückrichtung erfordert die verworfene Windungszahl. Der
Zeitpfeil ist keine Zusatzannahme, sondern Konsequenz der Geometrie des
Messens.

Das Pauli-Theorem gilt für einen offenen Zeitparameter. Für eine
kompakte Zeitdimension greift es nicht: Ihr Konjugiertes hat ein
diskretes Spektrum, die Windungszahlen. $\tilde T\cdot m = 1$
realisiert genau diese Situation: Zeit ist kompakt, das konjugierte
Spektrum ist die diskrete Massenleiter. Das Theorem schließt diesen Fall
nicht aus --- es erzwingt ihn als den einzigen konsistenten Weg, Zeit in
den Hilbertraum zu integrieren.

\section{Ausrollen: Typ I und Typ III}

Zwei verschiedene Operationen tragen denselben Namen. \textbf{Typ-I-
Ausrollen} ist der Messakt: verlustbehaftet, die Projektion vom
kompakten Träger auf die reelle Achse ist nicht umkehrbar. Das ist der
Kollaps. \textbf{Typ-III-Ausrollen} ist der Darstellungswechsel: der
bijektive Pullback aus Kapitel~\ref{ch:hilbert}, kein Informationsverlust,
nur ein Sprachwechsel. Der Korrekturfaktor $K_{\mathrm{frak}}$ sitzt in der
Massenformel des Typ-I-Übergangs; am Typ-III-Pullback ist er nicht
beteiligt. Wer in Hilbertraum-Sprache rechnet, sieht $K_{\mathrm{frak}}$
erst bei der SI-Übersetzung.\status{B}

\section{Zusammenfassung}

\begin{center}
\small\renewcommand{\arraystretch}{1.3}
\begin{tabular}{p{2.2cm}p{3.5cm}p{3.7cm}}
\toprule
 & \textbf{Standard-QM} & \textbf{FFGFT} \\
\midrule
Zeit & absent: kein Operator; externer Parameter & geometrisch: $\tilde T_i = 1/m_i$, kompakt, modengebunden \\
Superposition & zeitlos, keine intrinsische Uhr & intermediäre Torsion in einer Mode; $\tilde T_i$ als Skala \\
Modenübergreifend & algebraisch erlaubt, kein Zeitkonflikt & inkompatible $\tilde T_i\neq\tilde T_j$, strukturell problematisch \\
Verschränkung & zeitlos; Gleichzeitigkeit extern (Lorentz) & modenspezifisches $\tilde T$; kein gemeinsamer Zeitbegriff über Moden \\
Kollaps & instantan in $t$, keine Dauer & Typ I: gerichteter Informationsverlust, Zeitpfeil geometrisch \\
Pauli-Theorem & blockiert Zeitoperator für offenes $t$ & greift nicht für kompaktes $\tilde T$ \\
\bottomrule
\end{tabular}
\end{center}

Die Wahl, wo die Zeit sitzt, geht der Mathematik voraus und entscheidet,
ob die Moden geschenkt oder angeheftet sind. Wer die Sektoren als
räumliche Eigenräume liest, hält die vertraute Physik und zahlt mit
der äußeren Anheftung der Zahlen. Wer die Zeit als geschlossene
Dimension in den Raum nimmt, erhält die Moden strukturell und zahlt mit
einer ontologischen Festlegung. Der ontologische Preis der zweiten Wahl
bleibt auf dem Tisch --- aber der methodische Preis der ersten ist nicht
kleiner, er ist nur vertrauter.

\quelle{Dok.~307 (Zeit im Zustandsraum, Juli 2026), Dok.~334 (Superposition ohne Zeit, Aug.\ 2026), Dok.~333 (Typ I/III), Dok.~306 ($T\cdot E = 1$), Dok.~270 (Entkompaktifizierung)}
